% chapters/ch12_operational_refinements.tex
% Chapter 12: Operational Refinements + APCSO formal theory + applied guide
% Source: NRMO_v5_updated.pdf, Chapter 12 (lines 7331-7941 of v5_layout.txt)

\chapter{Operational Refinements --- APCSO Formal Theory and
Applied Guide}
\label{ch:operational-refinements}

This chapter completes the v2.1 operational refinements with three
additional sections:

\begin{itemize}
\item Limitations and future work (extending
Section~\ref{sec:invariants-limitations}).
\item APCSO --- A Chance-Constrained Autonomy-Preserving
Optimisation Framework for Non-Ruin Decision Shaping (formal
theory).
\item APCSO Operational \& Applied Guide --- from theory to
practice.
\end{itemize}

% =====================================================================
\section{Section 5.3: Abuse Detection Formalisation}
\label{sec:abuse-detection-formalisation}

\paragraph{Current state.}
Abuse patterns are detected via heuristics. No formal guarantee
exists that all misuse will be caught.

\paragraph{Known gaps.}
\begin{itemize}
\item Sophisticated misuse may evade pattern matching.
\item False positives require human review (resource cost).
\item Detection thresholds are implementation-dependent.
\end{itemize}

\paragraph{Future work.}
Develop formal verification methods for governance-layer integrity
through theoretical analysis and statistical anomaly detection.

% =====================================================================
\section{Section 5.4: Zero-Risk Classification}
\label{sec:zero-risk-classification}

\paragraph{Current state.}
Determining whether an action truly satisfies $\varepsilon = 0$ and
irreversibility $= 0$ requires human judgement. Edge cases exist.

\paragraph{Example ambiguity.}
A ``reversible'' social experiment may have irreversible cultural
effects. The boundary between zero-risk and low-risk can be
disputed.

\paragraph{Future work.}
Develop clearer taxonomies for action classification and establish
review processes for boundary cases.

% =====================================================================
\section{Section 5.5: Long-Term Co-Evolution}
\label{sec:long-term-coevolution}

\paragraph{Unsolved problem.}
NRMO and the institutional context co-evolve over decades. The
renewal protocol addresses parameter drift, but broader questions
remain:

\begin{itemize}
\item How does NRMO adapt to fundamental value shifts?
\item Can the system recognise when its own assumptions are
obsolete?
\item What triggers wholesale re-design vs incremental refinement?
\end{itemize}

\paragraph{Informal recommendation.}
Periodic reflective review (e.g., annual or decadal) asking:

\begin{quote}
\itshape
What did NRMO prevent that, in hindsight, should have been
permitted?
\end{quote}

\noindent
This is not formalised and relies on institutional discipline, not
algorithmic enforcement.

% =====================================================================
\section{Summary of Contributions (v2.1 Amendments)}
\label{sec:v21-summary}

NRMO v2.1 Amendments introduce four operational refinements to
address degradation patterns observed in long-term deployment:

\begin{enumerate}
\item \textbf{Parameter Renewal Protocol} prevents $\varepsilon$
ossification through periodic re-approval with mandatory log review.
\item \textbf{HOLD State Refinement} prevents indefinite stasis by
operationalising exit conditions.
\item \textbf{Zero-Risk Fast Path} enables computational optimisation
for trivially safe actions.
\item \textbf{Abuse Pattern Detection} provides operational
safeguards against vocabulary misuse.
\end{enumerate}

\subsection{Theoretical Integrity}

Core theoretical guarantees are preserved:

\begin{itemize}
\item Non-Ruin constraint
$P(\mathrm{Ruin} \mid \pi, \varepsilon) \le \varepsilon$ remains
enforced.
\item State space $S = \{\mathrm{ACTIVE}, \mathrm{HOLD},
\mathrm{EXIT}, \mathrm{SAFE\_EXIT}\}$ unchanged.
\item APCSO 3-choice output structure maintained.
\item Autonomy preservation principles intact.
\item No optimisation, ranking, or single-recommendation introduced.
\end{itemize}

\subsection{Operational Robustness}

Operational improvements achieved:

\begin{itemize}
\item Parameter evolution enabled while maintaining theoretical
bounds.
\item State stasis prevented through explicit exit conditions.
\item Creative actions facilitated without compromising safety.
\item Misuse patterns monitored and escalated.
\end{itemize}

\subsection{Design Philosophy}

The goal is not to make NRMO more permissive or more restrictive,
but to ensure it \emph{remains adaptive without compromising
invariants}. The system must preserve flexibility at the governance
layer while maintaining rigid constraints at the theoretical core.

\nrmobox{v2.1 design principle}{
\centering
\textit{The world can challenge NRMO's parameters, but not its
guarantees.}
}

\subsection{Future Directions}

Key areas for continued research and development:

\begin{itemize}
\item Formalisation of $\varepsilon$-governance mechanisms.
\item Verification methods for governance-layer integrity.
\item Automated detection of parameter drift and abuse patterns.
\item Long-term co-evolution frameworks for NRMO and institutional
contexts.
\item Empirical validation through extended deployment case
studies.
\end{itemize}

\subsection*{Document status}

\begin{tabular}{ll}
NRMO v2.1 Amendments Status: & Approved for Integration \\
Classification:               & Core Structure Corrections \\
Core Theory:                  & v2.0 Invariants Preserved \\
Recommended Adoption:         & Incremental deployment with
monitoring \\
\end{tabular}\\[0.4em]

\paragraph{Manifest.}
NRMO Core v1.2.7 / Passive Pattern v1.5 / APCSO (operational +
applied additions).
\textbf{Update Pack}: v1.2.8 (HTML operational changes).

% =====================================================================
\section{APCSO --- A Chance-Constrained, Autonomy-Preserving
Optimisation Framework for Non-Ruin Decision Shaping}
\label{sec:apcso-formal}

\nrmobox{Critical amendment: APCSO}{
\textbf{Structure}: APCSO operates between Passive Pattern and
NRMO Core.\\[0.3em]
Pipeline: PP $\to$ APCSO $\to$ NRMO. APCSO does \emph{not} replace
NRMO Core; it shapes the choice set before NRMO admissibility
verification.
}

\subsection{Abstract}

We introduce \emph{Autonomy-Preserving Choice-Set Optimisation}
(APCSO), a human-in-the-loop optimisation framework that shapes
decision \emph{environments} rather than selecting actions. APCSO
maximises mutual gain subject to explicit Non-Ruin constraints,
formally bounding the probability of irreversible breakdown. Human
autonomy is preserved by emitting unordered proposal sets instead of
ranked recommendations. We present a formal mathematical
formulation, complementary governance layers, a non-reducibility
argument, and an illustrative toy model.

\subsection{Introduction}

Many decision-support systems fail by over-optimising, collapsing
autonomy, or inducing irreversible commitment. APCSO addresses this
failure mode by optimising the \emph{structure} of available choices
under explicit ruin constraints. APCSO is architecturally separated
from suppressive safety layers and operates only where reversibility
remains.

\subsection{System Architecture}

\begin{itemize}
\item \textbf{NRMO Core}: Non-Ruin principle; no decision-making.
\item \textbf{Passive Pattern v1.5}: suppression-only safety layer.
\item \textbf{Autonomy-Preserving Choice-Set Optimisation (APCSO)}:
choice-set optimisation.
\item \textbf{$\varepsilon$-Governance / Policy Layer (external)}:
external authority for $\varepsilon$.
\item \textbf{Shadow Optimization Layer (non-executable)}.
\item \textbf{Autonomy Firewall / UI Layer}.
\item \textbf{Narrative / Ethics Layer (post-hoc)}.
\end{itemize}

\subsection{State, Observation, and Belief}

Let $s_t \in S$ denote a latent decision-relevant state. Observations
$o_t \in O$ induce a belief distribution:

\begin{equation}
b_t(s) = \Pr(s_t = s \mid o_{0:t}).
\label{eq:apcso-belief}
\end{equation}

\subsection{Proposal Set Structure}

APCSO outputs an unordered proposal set of exactly three elements:

\begin{equation}
\mathbf{a}_t = \{a_t^{(1)}, a_t^{(2)}, a_t^{(3)}\} \subset A.
\label{eq:apcso-proposal-set}
\end{equation}

\begin{itemize}
\item No ranking or preference is expressed.
\item At least one proposal must be a hold or exit option.
\end{itemize}

\subsection{Ruin Definition}

\begin{equation}
R = R_{\mathrm{irr}} \cup R_{\mathrm{aut}} \cup R_{\mathrm{lock}}.
\label{eq:apcso-ruin-set}
\end{equation}

Ruin is absorbing and irreversible once entered.

\subsection{Conservative Upper Bound on Ruin Probability}

\begin{equation}
\widehat{p}_R(a, b) =
\sup_{s \in \mathrm{supp}(b)} \Pr(R \mid s, a).
\label{eq:apcso-ruin-upper}
\end{equation}

\subsection{Mutual-Gain Utility}

\begin{equation}
U_{\mathrm{mg}}(s, a) =
\bigl(U_{\mathrm{self}}(s, a),\,
U_{\mathrm{other}}(s, a),\,
-I(a, s)\bigr).
\label{eq:apcso-utility}
\end{equation}

\noindent
Utility is evaluated only up to Pareto improvement.

\subsection{Optimisation Problem}

\begin{equation}
\max_{\mathbf{a} \subset A,\, |\mathbf{a}| = 3}
\mathbb{E}_{s \sim b}[U_{\mathrm{mg}}(s, \mathbf{a})]
\quad \text{s.t.} \quad
\widehat{p}_R(\mathbf{a}, b) \le \varepsilon.
\label{eq:apcso-optimisation}
\end{equation}

\subsection{Complementary Layers (Formal Theory)}

\subsubsection{$\varepsilon$-Governance / Policy Layer}

\begin{equation}
\varepsilon \in [\varepsilon_{\min}, \varepsilon_{\max}].
\end{equation}

Risk tolerance is determined externally by governance and never
optimised by APCSO.

\subsubsection{Shadow Optimisation}

\begin{equation}
a_t^{\mathrm{shadow}} =
\arg\max \mathbb{E}[U_{\mathrm{mg}}].
\end{equation}

Shadow results are never executable and are used solely for learning
and calibration.

\subsubsection{Autonomy Firewall / UI}

\begin{equation}
A_{\mathrm{afford}}(t) = \mathbf{a}_t.
\end{equation}

The firewall restricts availability without coercing choice.

\subsubsection{Narrative / Ethics Layer}

\begin{equation}
\frac{\partial \mathcal{N}}{\partial U_{\mathrm{mg}}} = 0,
\qquad
\frac{\partial \mathcal{N}}{\partial \varepsilon} = 0.
\end{equation}

Meaning attribution is strictly post-hoc and non-interfering.

\subsection{Non-Reducibility to Existing Frameworks}

APCSO is not reducible to CMDPs, Safe RL, or Mechanism Design. Those
frameworks optimise actions or policies; APCSO optimises the choice
set itself under autonomy constraints and outputs unordered feasible
sets.

\subsection{Toy Model Illustration}

Consider a two-state system $s \in \{\mathrm{safe},
\mathrm{fragile}\}$. An aggressive action is safe in isolation but
induces irreversible ruin when combined with commitment.
Action-centric optimisation selects it; APCSO excludes it by
structuring the choice set.

\subsection{Limitations}

\begin{itemize}
\item $\varepsilon$ is exogenous and governance-dependent.
\item Adversarial or self-destructive agents are out of scope.
\item Zero-sum or symbolic conflicts may not be expressible as
mutual-gain utilities.
\end{itemize}

\paragraph{Formal status.}
This document defines the complete Autonomy-Preserving Choice-Set
Optimisation (APCSO) theory, including complementary governance
layers, non-reducibility arguments, and an illustrative toy model.
NRMO Core and Passive Pattern v1.5 remain unchanged.

% =====================================================================
\section{APCSO Operational \& Applied Guide --- From Theory to
Practice}
\label{sec:apcso-applied-guide}

\subsection{Purpose of This Document}

This document operationalises the APCSO theory for real-world use.
It does not redefine the theory; instead, it specifies execution
patterns, governance hooks, and applied examples that can be safely
deployed without violating Non-Ruin guarantees.

\subsection{When APCSO Should Be Activated}

\begin{itemize}
\item Decisions with irreversible downside risk.
\item Human-in-the-loop judgement must be preserved.
\item Optimisation pressure exists but must be bounded.
\end{itemize}

\subsection{Operational Pipeline}

\begin{enumerate}
\item Passive Pattern Safety Check (halt if triggered).
\item Belief Update from observations.
\item $\varepsilon$ retrieval from Governance Layer.
\item Shadow Optimisation (optional, non-executable).
\item APCSO proposal-set optimisation.
\item Autonomy Firewall enforcement.
\item Narrative / Ethics explanation (post-hoc).
\end{enumerate}

\subsection{$\varepsilon$-Governance in Practice}

In deployment, $\varepsilon$ should be treated as a policy artefact,
not a tuning parameter. Typical stratification:

\begin{itemize}
\item Strategic decisions: $\varepsilon \in [0.01, 0.05]$.
\item Operational decisions: $\varepsilon \in [0.05, 0.15]$.
\item Exploratory / R\&D contexts: $\varepsilon \in [0.15, 0.30]$.
\end{itemize}

\subsection{Applied Examples (Industry-Specific)}

\subsubsection{Corporate Strategy Decision}

A board must choose whether to enter a new market with irreversible
capital commitments. APCSO structures proposals as:

\begin{itemize}
\item Market entry with capped exposure.
\item Pilot project with exit clause.
\item Hold / defer decision.
\end{itemize}

\subsubsection{AI System Deployment}

An AI team considers deploying a model with uncertain social impact.
APCSO ensures one proposal is a rollback-capable deployment and one
is non-deployment.

\subsubsection{Personal High-Stakes Decision}

In career or life decisions, APCSO prevents binary framing by forcing
a reversible alternative into the choice set.

\subsubsection{Public Policy / Regulation}

When introducing new regulations with irreversible social impact,
APCSO structures proposals as phased trials, sunset clauses, and
explicit non-action options, bounding systemic ruin.

\subsubsection{Healthcare / Clinical Governance}

For high-risk clinical or institutional decisions, APCSO enforces
reversible pilot protocols and ethics-review exit paths within the
proposal set.

\subsubsection{Finance / Capital Allocation}

In capital deployment, APCSO ensures that at least one proposal
preserves liquidity and prevents lock-in, mitigating tail-risk
exposure.

\subsection{Metrics for Ongoing Monitoring}

\begin{itemize}
\item Observed vs estimated ruin probability.
\item Frequency of hold / exit selections.
\item Shadow vs executable divergence.
\end{itemize}

\subsection{Failure Modes and Safeguards}

\begin{itemize}
\item $\varepsilon$ drift without governance review.
\item Shadow optimisation leaking into execution.
\item UI bypass of Autonomy Firewall.
\end{itemize}

\subsection{Operational Checklist (SOP)}

\begin{enumerate}
\item Confirm APCSO applicability (irreversibility and ruin risk
present).
\item Run Passive Pattern safety check (halt if triggered).
\item Retrieve $\varepsilon$ from Governance Layer (document
authority).
\item Execute Shadow Optimisation strictly offline.
\item Generate exactly three proposals (incl. hold / exit).
\item Validate ruin upper-bound $\le \varepsilon$ for all proposals.
\item Enforce Autonomy Firewall at UI / process level.
\item Log human choice without preference inference.
\item Attach Narrative explanation post-hoc only.
\end{enumerate}

\subsection{End-to-End Operational Scenario (90-Day Example)}

\paragraph{Scenario.}
New business initiative with irreversible capital exposure.

\begin{description}
\item[Day 0] APCSO activation approved; $\varepsilon$ set by
executive governance.
\item[Day 7] Shadow Optimisation explores aggressive expansion
(non-executable).
\item[Day 14] APCSO generates proposal set: capped entry / pilot /
hold.
\item[Day 30] Human selects pilot option; monitoring begins.
\item[Day 60] Ruin indicators reviewed; $\varepsilon$ unchanged.
\item[Day 90] Exit clause exercised or scaled entry proposed.
\end{description}

\subsection{Failure Scenarios (Negative Cases)}

\subsubsection{$\varepsilon$ Drift Failure}

$\varepsilon$ is gradually increased without governance review,
leading to silent accumulation of ruin risk. This violates
$\varepsilon$-Governance separation and invalidates APCSO
guarantees.

\subsubsection{Shadow Leakage Failure}

Shadow Optimisation outputs are exposed to decision-makers, biasing
human choice toward aggressive actions. This collapses autonomy and
reintroduces covert ranking.

\subsubsection{Hold-Option Erosion}

Hold / exit proposals are repeatedly framed as inferior, effectively
removing reversibility. APCSO degenerates into action optimisation.

\subsection{Role and Responsibility Decomposition}

\begin{itemize}
\item \textbf{Governance Authority}: defines and revises
$\varepsilon$.
\item \textbf{APCSO Operator}: generates proposal sets and validates
constraints.
\item \textbf{Shadow Analyst}: runs non-executable exploration only.
\item \textbf{Human Decision-Maker}: selects freely among proposals.
\item \textbf{Auditor}: verifies Autonomy Firewall and process
compliance.
\item \textbf{Narrative Owner}: provides post-hoc explanation only.
\end{itemize}

\subsection{Advanced Applications (Applied Extensions)}

\subsubsection{Multi-Stage APCSO (Recursive Deployment)}

Multi-Stage APCSO addresses decisions whose consequences unfold over
extended horizons. Let the outcome of stage $k$ update the belief
$b_{k+1}$:

\begin{equation}
b_{k+1} = \mathrm{Update}(b_k, o_k).
\end{equation}

APCSO is re-invoked at each stage. The key constraint is that
$\varepsilon$ remains fixed across stages unless explicitly updated
by Governance. This prevents gradual normalisation of risk. Recursive
APCSO therefore enables acceleration while preserving long-term
Non-Ruin guarantees.

\subsubsection{Portfolio-Level APCSO (Ruin Budget Allocation)}

Organisations often face multiple concurrent decisions.
Portfolio-Level APCSO introduces a global ruin budget
$\varepsilon_{\mathrm{total}}$, allocated across $N$ initiatives:

\begin{equation}
\sum_{i=1}^N \varepsilon_i \le \varepsilon_{\mathrm{total}}.
\end{equation}

Each local APCSO instance must satisfy its assigned $\varepsilon_i$.
This enforces systemic stability even when individual projects
pursue aggressive trajectories.

\subsubsection{Crisis Mode Operation}

In crisis conditions (time pressure, incomplete information), APCSO
operates in a constrained mode. Proposal diversity may be reduced,
but invariants must hold:

\begin{itemize}
\item At least one explicit hold / exit option.
\item No $\varepsilon$ expansion; $\varepsilon$ tightening is
preferred.
\item Shadow Optimisation remains isolated.
\end{itemize}

\subsubsection{Human Bias Mitigation Layer}

APCSO should be paired with bias-mitigation techniques: randomised
order, neutral framing, and visual symmetry.

\subsubsection{Institutional Scaling and Federation}

Local APCSO units operate under a global $\varepsilon$ ceiling
enforced by central governance:

\begin{equation}
\varepsilon_{\mathrm{local}} \le \varepsilon_{\mathrm{global}}.
\end{equation}

\subsubsection{Adversarial and Political Environments}

APCSO may be applied internally in adversarial settings to structure
irreversible commitments while treating adversaries as uncertainty.

\subsubsection{Termination and Exit Protocols}

If belief uncertainty exceeds model validity or $\varepsilon$ cannot
be supplied, APCSO must suspend and defer to Passive Pattern halt.

\subsection{Implementation Mapping (Code Correspondence Table)}

\begin{table}[ht]
\centering
\small
\begin{tabularx}{\linewidth}{lXX}
\toprule
\textbf{Theory component} & \textbf{Implementation layer} & \textbf{Typical artefacts} \\
\midrule
Belief State $b_t$
& State Estimation
& Bayesian updater, belief store, uncertainty tracker \\

$\varepsilon$-Governance
& Policy / Config
& Policy file, approval workflow, configuration service \\

Shadow Optimisation
& Offline Analytics
& Simulation jobs, notebooks, sandboxed optimisers \\

APCSO Optimiser
& Core Service
& Proposal generator, constraint checker \\

Ruin Upper Bound
& Risk Engine
& Worst-case estimator, safety validator \\

Autonomy Firewall
& Interface / Process
& UI gating, API allow-list, procedural controls \\

Narrative Layer
& Explanation / Reporting
& Post-hoc report generator, audit logs \\

Passive Pattern Safety
& Interrupt
& Kill-switch, circuit breaker \\

\bottomrule
\end{tabularx}
\caption[APCSO implementation mapping]{APCSO theoretical components
and their typical implementation artefacts. The mapping is
illustrative and does not mandate a specific technology stack.}
\label{tab:apcso-implementation}
\end{table}

This mapping intentionally avoids executable pseudo-code. Its
purpose is to clarify separation-of-concerns and prevent unsafe
coupling between layers.

\subsection{Product Archetype and Design Boundaries}

This section defines the abstract product design compatible with
APCSO. It intentionally avoids concrete specifications, UI mockups,
or implementation details. Its purpose is to prevent misconstruction
while enabling safe instantiation.

\subsubsection{Product Archetype}

An APCSO-compatible product is a Decision Support Infrastructure,
not a decision-maker. Its core function is to structure feasible
choices under Non-Ruin constraints, while explicitly delegating
judgement to humans or institutions.

\begin{itemize}
\item Primary role: choice-set generation and validation.
\item Non-role: action selection, ranking, or execution.
\end{itemize}

\subsubsection{Acceptable Product Forms}

\begin{itemize}
\item Internal governance tool.
\item Advisory layer embedded in existing workflows.
\item Audit and review infrastructure.
\end{itemize}

\subsubsection{Prohibited Product Patterns}

\begin{itemize}
\item Automatic execution engines.
\item Recommendation systems with ranked outputs.
\item Interfaces that expose $\varepsilon$ tuning to end-users.
\item Dashboards that surface Shadow Optimisation results.
\end{itemize}

\subsubsection{Human-Product Boundary}

The product must preserve a clear separation:

\begin{itemize}
\item The system structures options.
\item The human selects and bears responsibility.
\end{itemize}

\subsubsection{Governance and Accountability}

The product must make visible who supplied $\varepsilon$, who
approved APCSO activation, and when termination conditions apply.
However, it must not enforce or suggest normative judgements.

\subsubsection{Anti-Pattern Warning}

If a product design attempts to optimise user behaviour, persuasion,
or compliance, it violates APCSO's autonomy-preserving premise and
must be rejected.

\subsection{Relationship to Core Theory}

All operational practices here respect the formal constraints of the
APCSO theory and must be suspended if Passive Pattern v1.5 or NRMO
Core triggers a halt. This operational guide is a companion document
to the formal APCSO theory. It is non-normative, non-binding, and
intended solely for safe application.

% =====================================================================
\section*{Connections to Other Parts}

\begin{itemize}
\item APCSO is the formal mathematical core of the choice-set system
referenced in the v2.1 invariants
(Section~\ref{sec:inv-apcso}).
\item Equation~\ref{eq:apcso-optimisation} (the constrained
optimisation problem) is the action-side analogue of the
admissibility-set construction
$\Pi_{\mathrm{safe}}(T, \varepsilon)$ in
Section~\ref{sec:foundations-3-1-definition}.
\item The chance constraint
$\widehat{p}_R(\mathbf{a}, b) \le \varepsilon$ in
Equation~\ref{eq:apcso-ruin-upper} is the local form of the
robust chance constraint $\sup_{P \in \mathcal{P}}
P_{\pi, P}(\tau_F \le T) \le \varepsilon$ used in the arXiv paper
(Part~VII).
\item The complementary governance layers (Shadow Optimisation,
Autonomy Firewall, Narrative Layer) correspond to the
governance--execution separation invariant
(Invariant~\ref{inv:gov-exec-sep}).
\item The Multi-Stage APCSO recursion is the cognitive analogue of
the Long-Horizon Drift Control of $\Omega$ Full v5.0 (Part~IX),
which uses a sustainable-growth estimator and
$\lambda_{\mathrm{drift}}$ penalty to prevent gradual normalisation
in extended horizons.
\end{itemize}
